%% bab4.tex
%% BAB 4 - HASIL DAN PEMBAHASAN
%% Format sesuai Panduan Penyusunan Tesis ITS 2021

\chapter{HASIL DAN PEMBAHASAN}

Pada awal bagian ini, perlu diberikan suatu pengantar yang memuat hal-hal yang akan dilakukan beserta analisis yang digunakan dalam menyelesaikan penelitian. Selanjutnya secara terperinci dan tahap demi tahap tujuan penelitian dibahas dan dianalisis secara detail dan tajam, dengan menggunakan metode yang telah diberikan dalam metodologi penelitian, sampai diperoleh suatu hasil penelitian.

\section{Hasil Pengumpulan Data}

Jelaskan hasil pengumpulan data yang telah dilakukan. Sajikan data dalam bentuk tabel atau grafik yang jelas dan mudah dipahami.

\subsection{Data Pertama}

Tabel~\ref{tab:data-pertama} menunjukkan hasil pengumpulan data pertama.

\begin{table}[htbp]
    \caption{Hasil Pengumpulan Data Pertama}
    \label{tab:data-pertama}
    \centering
    \begin{tabular}{|c|c|c|c|}
        \hline
        \textbf{No} & \textbf{Parameter} & \textbf{Nilai} & \textbf{Satuan} \\
        \hline
        1 & Parameter A & 100 & Unit A \\
        \hline
        2 & Parameter B & 200 & Unit B \\
        \hline
        3 & Parameter C & 300 & Unit C \\
        \hline
    \end{tabular}
    
    \vspace{0.3cm}
    \footnotesize{Sumber: Hasil Pengamatan}
\end{table}

\subsection{Data Kedua}

Jelaskan data kedua yang telah dikumpulkan.

\section{Implementasi}

Jelaskan implementasi dari metode yang digunakan dalam penelitian ini.

\subsection{Implementasi Tahap Pertama}

Jelaskan implementasi tahap pertama.

\subsection{Implementasi Tahap Kedua}

Jelaskan implementasi tahap kedua.

\section{Hasil Pengujian}

Jelaskan hasil pengujian yang telah dilakukan.

\subsection{Skenario Pengujian}

Jelaskan skenario pengujian yang digunakan.

\subsection{Hasil Pengujian}

Gambar~\ref{fig:hasil-pengujian} menunjukkan hasil pengujian yang telah dilakukan.

\begin{figure}[htbp]
    \centering
    % Ganti dengan gambar hasil pengujian Anda
    \fbox{\parbox{0.8\textwidth}{\centering\vspace{3cm}[Grafik Hasil Pengujian]\vspace{3cm}}}
    \caption{Hasil Pengujian}
    \label{fig:hasil-pengujian}
\end{figure}

\section{Pembahasan}

Analisis dan pembahasan ini dilakukan untuk semua tujuan yang telah ditetapkan pada tujuan penelitian.

\subsection{Analisis Hasil Pertama}

Jelaskan analisis terhadap hasil pertama. Bandingkan dengan penelitian sebelumnya dan jelaskan temuan yang diperoleh.

\subsection{Analisis Hasil Kedua}

Jelaskan analisis terhadap hasil kedua.

\subsection{Analisis Hasil Ketiga}

Jelaskan analisis terhadap hasil ketiga.

\section{Evaluasi}

Jelaskan evaluasi keseluruhan dari hasil penelitian ini. Bandingkan dengan target yang telah ditetapkan di awal penelitian.
